\documentclass[tikz,border=10pt]{standalone}
\usepackage{tikz-3dplot}

\begin{document}
\tdplotsetmaincoords{70}{110}
\def\l{2.0} 
\def\radius{8pt}
\def\scal{4.0}
\def\colorconnectivity{blue} 

\input{params.tex}
\begin{tikzpicture}[font=\LARGE,tdplot_main_coords, every node/.style={scale=\scal}]

\def\locations{1,...,20}

\coordinate (1) at (0,-2\l,0);
\coordinate (2) at (-2\l,-2\l,0);
\coordinate (3) at (-2\l,0,0);
\coordinate (4) at (0,0,0);

\coordinate (5) at (0,-2\l,2\l);
\coordinate (6) at (-2\l,-2\l,2\l);
\coordinate (7) at (-2\l,0,2\l);
\coordinate (8) at (0,0,2\l);

\coordinate (9)  at ($(1)!0.5!(2)$);
\coordinate (10) at ($(2)!0.5!(3)$);
\coordinate (11) at ($(3)!0.5!(4)$);
\coordinate (12) at ($(4)!0.5!(1)$);
\coordinate (13) at ($(5)!0.5!(6)$);
\coordinate (14) at ($(6)!0.5!(7)$);
\coordinate (15) at ($(7)!0.5!(8)$);
\coordinate (16) at ($(8)!0.5!(5)$);
\coordinate (17) at ($(1)!0.5!(5)$);
\coordinate (18) at ($(2)!0.5!(6)$);
\coordinate (19) at ($(3)!0.5!(7)$);
\coordinate (20) at ($(8)!0.5!(4)$);


\draw[-]
(3) -- (4)
(3) -- (7)
(4) -- (1)
(5) -- (6)
(6) -- (7)
(7) -- (8)
(8) -- (5)
(1) -- (5)
(4) -- (8);

\draw[-, dashed]
(1) -- (2)
(2) -- (3)
(2) -- (6);


\foreach \x [count=\c,evaluate=\c as \y using {{\connectivityHexaQ}[\c-1]}]  in \locations 
         {
           \fill[black] (\x) circle [radius=\radius];
           \node[\colorconnectivity, above left] at (\x) {\y};
         }

\end{tikzpicture}
\end{document}
